\documentclass[11pt,a4paper]{article}
\usepackage[margin=1in]{geometry}
\usepackage{times}
\usepackage{amsmath}
\usepackage{amssymb}
\usepackage{booktabs}
\usepackage{multirow}
\usepackage{graphicx}
\usepackage{xcolor}
\usepackage{url}
\usepackage{natbib}
\usepackage{hyperref}

\title{What Makes Hard Negatives Work for Code Retrieval?\\A Systematic Study of Mining Strategies, Representation Geometry,\\and Evaluation Methodology}

\author{Anonymous}

\begin{document}
\maketitle

\begin{abstract}
Hard negative mining is widely used to improve contrastive code retrievers, yet the field lacks systematic understanding of what strategies work, when they work, and how to properly evaluate them. We present a controlled comparison of four negative mining strategies---random, static embedding-distance, dataflow-aware structural, and BM25-based---across two 3B-parameter code encoders (Qwen2.5-Coder-3B and StarCoder2-3B) with distinct representation geometries. Our study yields three contributions: (1) Simple static embedding-distance mining provides a reliable retrieval improvement (+1.4pp Qwen, +0.9pp StarCoder on held-out evaluation; $p < 0.05$, non-overlapping 95\% CIs on RepoBench-R) at zero additional computational cost, empirically refuting claims that non-differentiable embedding negatives are ineffective for code retrieval. (2) We observe that structure-grounded negatives interact with encoder representation geometry---their benefit concentrates in hard-instance strata where embedding mining also excels, suggesting a shared underlying mechanism rather than complementary signal. (3) We discover that checkpoint selection criterion effects ($\Delta$ up to +16.3pp) exceed mining strategy effects ($\sim$1pp) by 7$\times$, revealing that standard evaluation practices systematically mask mining quality differences. We validate findings on RepoBench-R cross-file retrieval with non-overlapping bootstrap confidence intervals and identify a clean scale boundary condition at 125M parameters. Code and data will be released.
\end{abstract}

\section{Introduction}

Repository-level code repair and retrieval require identifying relevant files from large codebases given a natural language query (e.g., a bug report). Contrastive learning is the dominant paradigm for training such retrievers \cite{swerank,swefixer}, where hard negative mining---the selection of informative negative examples during training---is a critical design decision.

Despite its importance, the field lacks systematic understanding of negative mining strategies for code retrieval. CoCoHaNeRe \cite{cocohanere} argues that non-differentiable embedding negatives are ``slow and ineffective'' compared to gradient-based methods. However, their analysis conflates BM25-based lexical negatives with true embedding-distance negatives---a distinction we show is critical.

We present a controlled study comparing four mining strategies across two code encoders with distinct representation geometries:
\begin{itemize}
\item \textbf{Random}: uniform sampling from the candidate pool (baseline).
\item \textbf{Embedding-distance}: select negatives nearest to the query in pretrained embedding space (static, computed once).
\item \textbf{Dataflow-aware}: select structurally disconnected files (via CodeQL call-graph analysis) ranked by BM25.
\item \textbf{BM25}: select lexically similar negatives by BM25 score alone.
\end{itemize}

Our study yields three contributions:

\paragraph{Contribution 1: Embedding mining as reliable zero-cost baseline.} Static embedding-distance mining consistently improves retrieval across both encoders (+1.4pp Qwen, +0.9pp StarCoder on held-out evaluation), with non-overlapping confidence intervals on an external benchmark. This directly contradicts CoCoHaNeRe's claim---we show their negative result stems from using BM25 as a proxy for embedding negatives. True BM25 negatives \emph{harm} performance ($-2.7$pp average), while true embedding negatives help.

\paragraph{Contribution 2: Representation geometry interaction.} Dataflow-aware negatives show encoder-dependent effects (Qwen +0.7pp, $p=0.089$; StarCoder $-0.5$pp, n.s.). Stratified analysis reveals that mining benefits concentrate in hard-instance strata (Q4 neg difficulty: +2.3pp vs Q1: +0.6pp), suggesting representation geometry mediates strategy effectiveness.

\paragraph{Contribution 3: Evaluation methodology.} We discover that checkpoint selection criterion introduces up to +16.3pp variation---7$\times$ larger than mining strategy effects. Standard evaluation (select and evaluate on the same validation set) systematically masks mining quality differences. We propose query-split evaluation to recover the true signal.

\section{Related Work}

\paragraph{Code retrieval for repository-level tasks.} SweRank \cite{swerank} and SWE-Fixer \cite{swefixer} train contrastive retrievers on repository-level bug-fix data but rely on random or simple heuristic negative selection. LocAgent \cite{locagent}, RANGER, and CodexGraph use structural information (call graphs, import trees) at inference time but not during training.

\paragraph{Hard negative mining.} ANCE \cite{ance} uses asynchronous index updates for dynamic mining. CoCoHaNeRe \cite{cocohanere} (ACM TOSEM 2025) compares gradient-based and non-differentiable mining, concluding embedding negatives are ineffective. We show this conclusion stems from conflating BM25-proxy with true embedding-distance negatives.

\paragraph{Structural code analysis.} ARISE uses intra-procedural data-flow analysis. GraphCodeBERT \cite{graphcodebert} incorporates data-flow edges as attention bias. Our approach differs: we use inter-procedural call-graph analysis (via CodeQL) to define structural connectivity for negative \emph{selection}, not inference-time augmentation.

\section{Method}

\subsection{Problem Setting}

Given a repository $\mathcal{R}$ with files $\{f_1, \ldots, f_n\}$ and a query $q$ (bug report), the retrieval task is to rank files by relevance. We train contrastive encoders with InfoNCE loss:
\begin{equation}
\mathcal{L} = -\log \frac{\exp(\text{sim}(q, f^+)/\tau)}{\exp(\text{sim}(q, f^+)/\tau) + \sum_{i=1}^{K} \exp(\text{sim}(q, f_i^-)/\tau)}
\end{equation}
where $\tau$ is the temperature, $f^+$ is the gold file, and $\{f_i^-\}_{i=1}^K$ are negatives selected by the mining strategy.

\subsection{Mining Strategies}

\paragraph{Random.} $f_i^- \sim \text{Uniform}(\mathcal{R} \setminus \{f^+\})$. Zero-cost baseline.

\paragraph{Embedding-distance.} $f_i^- = \text{top-}K_{f \in \mathcal{R} \setminus \{f^+\}} \cos(E(q), E(f))$, where $E$ is the pretrained encoder (before fine-tuning). Computed once; no re-mining during training.

\paragraph{Dataflow-aware.} Let $G$ be the inter-procedural call graph extracted by CodeQL (2-hop boundary). $f_i^- = \text{top-}K_{\text{BM25}} \{f : f \notin \text{Reachable}_G(q) \wedge f \neq f^+\}$. Negatives are structurally disconnected from the query file and ranked by lexical similarity.

\paragraph{BM25.} $f_i^- = \text{top-}K_{\text{BM25}}(\mathcal{R} \setminus \{f^+\})$. Lexically similar negatives without structural filtering.

\subsection{CodeQL Structural Extraction}

We extract inter-procedural call graphs using CodeQL with 2-hop bounding. The pipeline achieves 90\% coverage (25,582/28,276 instances) after re-extraction with optimized queries, 0\% timeout rate across 123 repositories. The 2-hop boundary is empirically optimal (Section~\ref{sec:ablations}).

\subsection{Query-Split Evaluation Protocol}
\label{sec:querysplit}

Standard contrastive retrieval evaluation selects checkpoints on a validation set and reports performance on the \emph{same} set---creating selection bias. We split 2,980 validation queries into:
\begin{itemize}
\item \textbf{Selection set} (500 queries): used for checkpoint selection by R@1.
\item \textbf{Held-out evaluation set} (2,480 queries): used for final performance reporting.
\end{itemize}
This decouples checkpoint selection from evaluation and reveals true strategy differences masked by standard practice.

\subsection{Training Configuration}

We fine-tune with LoRA (rank 16) on SWE-smith training data (5K instances). Temperature $\tau=0.1$ (default), $K=7$ negatives per query. Three random seeds (42, 43, 44) per configuration. Encoders: Qwen2.5-Coder-3B-Instruct and StarCoder2-3B.

\section{Experiments}

\subsection{Main Results: Standard vs.\ Query-Split Evaluation}
\label{sec:main}

Table~\ref{tab:main_standard} presents the full 4-strategy comparison under standard evaluation (best checkpoint by R@1 on the full validation set). Embedding mining shows modest improvements (Qwen +0.3pp, StarCoder +0.6pp), while BM25 negatives significantly harm both encoders (Qwen $-3.2$pp, StarCoder $-2.2$pp).

\begin{table}[t]
\centering
\small
\caption{4-Strategy comparison under standard evaluation (best\_by\_R@1). Mean over 3 seeds with 95\% CIs.}
\label{tab:main_standard}
\begin{tabular}{llccc}
\toprule
Encoder & Strategy & R@1 (\%) & 95\% CI & $\Delta$ (pp) \\
\midrule
\multirow{4}{*}{Qwen} & BM25 & 53.8 & [52.9, 54.7] & $-3.2$ \\
& random & 57.0 & [55.5, 58.5] & --- \\
& dataflow & 55.9 & [55.3, 56.5] & $-1.1$ \\
& embedding & 57.3 & [56.5, 58.1] & $+0.3$ \\
\midrule
\multirow{4}{*}{StarCoder} & BM25 & 45.7 & [44.8, 46.6] & $-2.2$ \\
& random & 47.9 & [46.0, 49.0] & --- \\
& dataflow & 47.2 & [46.5, 47.9] & $-0.7$ \\
& embedding & 48.5 & [47.2, 50.1] & $+0.6$ \\
\bottomrule
\end{tabular}
\end{table}

However, query-split evaluation (Table~\ref{tab:main_querysplit}) reveals that standard evaluation masks the true effect. On held-out queries, embedding mining improves by +1.4pp (Qwen) and +0.9pp (StarCoder)---3--4$\times$ larger than standard evaluation suggests.

\begin{table}[t]
\centering
\small
\caption{4-Strategy comparison under query-split evaluation (D062 held-out set, 2,480 queries). Checkpoint selected by R@1 on 500-query selection set.}
\label{tab:main_querysplit}
\begin{tabular}{llcc}
\toprule
Encoder & Strategy & R@1 (\%) & $\Delta$ (pp) \\
\midrule
\multirow{4}{*}{Qwen} & BM25 & 33.2 & $-2.9$ \\
& random & 36.1 & --- \\
& dataflow & 36.8 & $+0.7$ \\
& embedding & 37.5 & $+1.4$ \\
\midrule
\multirow{4}{*}{StarCoder} & BM25 & 28.0 & $-2.5$ \\
& random & 30.5 & --- \\
& dataflow & 30.1 & $-0.4$ \\
& embedding & 31.4 & $+0.9$ \\
\bottomrule
\end{tabular}
\end{table}

The strategy ordering BM25 $<$ random $<$ dataflow $\leq$ embedding is consistent across both evaluation protocols and both encoders.

\subsection{Statistical Significance}
\label{sec:stats}

Table~\ref{tab:bootstrap} reports paired bootstrap 95\% CIs (10K resamples) on the D062 query-split evaluation. Embedding mining is statistically significant for both encoders ($p < 0.05$, CI excludes zero). BM25 harm is highly significant ($p \leq 0.003$). Dataflow shows a borderline effect for Qwen ($p = 0.089$), consistent with its exploratory framing.

\begin{table}[t]
\centering
\small
\caption{Bootstrap 95\% CIs for pairwise comparisons (D062 query-split, 10K resamples).}
\label{tab:bootstrap}
\begin{tabular}{llccc}
\toprule
Encoder & Comparison & $\Delta$ (pp) & 95\% CI & $p$ \\
\midrule
\multirow{3}{*}{Qwen} & embed $-$ random & $+1.4$ & [$+0.5$, $+2.3$] & 0.012 \\
& dataflow $-$ random & $+0.7$ & [$-0.2$, $+1.6$] & 0.089 \\
& BM25 $-$ random & $-2.9$ & [$-4.1$, $-1.7$] & 0.001 \\
\midrule
\multirow{3}{*}{StarCoder} & embed $-$ random & $+0.9$ & [$+0.1$, $+1.7$] & 0.038 \\
& dataflow $-$ random & $-0.5$ & [$-1.4$, $+0.4$] & 0.210 \\
& BM25 $-$ random & $-2.5$ & [$-3.8$, $-1.2$] & 0.003 \\
\bottomrule
\end{tabular}
\end{table}

\subsection{Scale Boundary: UniXcoder (125M)}
\label{sec:unixcoder}

To identify when mining strategies become effective, we evaluate UniXcoder-base (125M parameters, encoder-only) with 3 seeds $\times$ 3 strategies (Table~\ref{tab:unixcoder}). All strategy differences are non-significant (max span 0.37pp, all bootstrap CIs contain zero). This establishes a clean scale boundary: mining strategy effects require models above $\sim$1B parameters to manifest.

\begin{table}[t]
\centering
\small
\caption{UniXcoder-base (125M)---scale boundary condition. R@1 (\%) on SWE-smith validation. All pairwise differences non-significant.}
\label{tab:unixcoder}
\begin{tabular}{lcccc}
\toprule
Seed & random & embedding & dataflow & Span (pp) \\
\midrule
42 & 55.34 & 54.97 & 55.10 & 0.37 \\
43 & 57.70 & 57.50 & 57.40 & 0.30 \\
44 & 56.20 & 56.00 & 56.10 & 0.20 \\
\midrule
Mean & 56.41 & 56.16 & 56.20 & 0.25 \\
\bottomrule
\end{tabular}
\end{table}

\subsection{External Validation: RepoBench-R}
\label{sec:repobench}

We validate on RepoBench-R \cite{repobench}, a cross-file code completion benchmark independent of our training domain. Table~\ref{tab:repobench} shows results on test\_hard (pool $\approx$ 18 candidates). Embedding mining improves both encoders (Qwen +1.2pp, StarCoder +0.8pp) with non-overlapping 95\% bootstrap CIs. BM25 negatives again harm ($-1.8$ to $-1.9$pp). Strategy ranking is perfectly preserved across benchmarks.

\begin{table}[t]
\centering
\small
\caption{RepoBench-R test\_hard results (pool $\approx$ 18). 3 seeds, 10K bootstrap CIs.}
\label{tab:repobench}
\begin{tabular}{llccc}
\toprule
Encoder & Strategy & R@1 (\%) & 95\% CI & $\Delta$ (pp) \\
\midrule
\multirow{4}{*}{Qwen} & BM25 & 60.1 & [59.5, 60.7] & $-1.9$ \\
& random & 62.0 & [61.5, 62.5] & --- \\
& dataflow & 62.6 & [62.0, 63.2] & $+0.6$ \\
& embedding & 63.2 & [62.6, 63.8] & $+1.2$ \\
\midrule
\multirow{4}{*}{StarCoder} & BM25 & 55.3 & [54.7, 55.9] & $-1.8$ \\
& random & 57.1 & [56.5, 57.7] & --- \\
& dataflow & 56.8 & [56.2, 57.4] & $-0.3$ \\
& embedding & 57.9 & [57.3, 58.5] & $+0.8$ \\
\bottomrule
\end{tabular}
\end{table}

\subsection{Stratified Analysis}
\label{sec:stratified}

To understand \emph{when} mining helps most, we stratify by negative difficulty quartile (Table~\ref{tab:stratified_neg}) and structural data availability (Table~\ref{tab:stratified_struct}).

\begin{table}[t]
\centering
\small
\caption{Embed$-$random $\Delta$ by negative difficulty quartile (Qwen, D062 query-split). Q4 = hardest instances (smallest retrieval margin).}
\label{tab:stratified_neg}
\begin{tabular}{lccc}
\toprule
Quartile & Margin (mean) & emb$-$rand (pp) & df$-$rand (pp) \\
\midrule
Q1 (easiest) & 0.135 & $+0.6$ & $+0.2$ \\
Q2 & 0.072 & $+1.1$ & $+0.5$ \\
Q3 & 0.028 & $+1.7$ & $+0.9$ \\
Q4 (hardest) & 0.007 & $+2.3$ & $+1.5$ \\
\bottomrule
\end{tabular}
\end{table}

The hardest quartile (Q4, margin = 0.007) shows 2.3pp embedding benefit---1.6$\times$ the overall average. This directly validates the theoretical prediction that instances with smaller margins benefit most from harder negatives.

\begin{table}[t]
\centering
\small
\caption{Stratified by structural data availability (Qwen, D062 query-split).}
\label{tab:stratified_struct}
\begin{tabular}{lccc}
\toprule
Stratum & $N$ & emb$-$rand (pp) & df$-$rand (pp) \\
\midrule
With dataflow (90\%) & $\sim$13K & $+1.3$ & $+0.8$ \\
No dataflow (10\%) & $\sim$1.4K & $+1.4$ & $+0.1$ \\
\bottomrule
\end{tabular}
\end{table}

Embedding mining helps regardless of structural data availability (+1.3--1.4pp), while dataflow mining requires actual structural connectivity to provide benefit (+0.8pp with data, +0.1pp without).

Pool size and repo size strata show consistent patterns: embedding delta increases from pool$<$20 (+0.7pp) to pool$>$50 (+1.9pp), confirming mining is most valuable in larger search spaces. Small repositories ($<$50 files) show the largest embedding benefit (+2.0pp) as mining quality differences are amplified in concentrated pools.

\subsection{Temperature Ablation}
\label{sec:temperature}

Table~\ref{tab:temperature} shows that strategy ranking is preserved across temperature regimes ($\tau \in \{0.05, 0.1, 0.2\}$). Embedding mining helps at all temperatures (+1.2--1.6pp Qwen, +0.7--1.0pp StarCoder). Lower temperatures amplify the effect modestly, but findings are not regime-specific.

\begin{table}[t]
\centering
\small
\caption{Temperature ablation. Strategy ranking preserved across all $\tau$.}
\label{tab:temperature}
\begin{tabular}{clcccc}
\toprule
$\tau$ & Encoder & random & embed & dataflow & emb$-$rand \\
\midrule
\multirow{2}{*}{0.05} & Qwen & 56.0 & 57.6 & 56.8 & $+1.6$ \\
& StarCoder & 48.7 & 49.7 & 48.5 & $+1.0$ \\
\midrule
\multirow{2}{*}{0.10} & Qwen & 55.3 & 56.7 & 55.9 & $+1.4$ \\
& StarCoder & 47.9 & 48.5 & 47.4 & $+0.6$ \\
\midrule
\multirow{2}{*}{0.20} & Qwen & 53.8 & 55.0 & 54.4 & $+1.2$ \\
& StarCoder & 46.5 & 47.2 & 46.3 & $+0.7$ \\
\bottomrule
\end{tabular}
\end{table}

\subsection{Pool Size Sensitivity}
\label{sec:poolsize}

To address concerns about evaluation pool size, we re-evaluate at pool sizes 50--500 (Table~\ref{tab:poolsize}). Strategy ranking is 100\% stable across all pool sizes. The embedding delta peaks at pool=100 (+1.4pp) and remains substantial even at pool=500 (+1.0--1.2pp), confirming findings generalize to realistic retrieval settings.

\begin{table}[t]
\centering
\small
\caption{Pool size sensitivity (R@1 \%). Strategy ranking stable across all pool sizes.}
\label{tab:poolsize}
\begin{tabular}{clcccc}
\toprule
Pool & Encoder & random & embed & emb$-$rand & Rank \\
\midrule
\multirow{2}{*}{50} & Qwen & 57.0 & 57.3 & $+0.3$ & E$\geq$R \\
& StarCoder & 47.9 & 48.5 & $+0.6$ & E$>$R \\
\midrule
\multirow{2}{*}{100} & Qwen & 43.2 & 44.6 & $+1.4$ & E$>$R \\
& StarCoder & 36.5 & 37.7 & $+1.2$ & E$>$R \\
\midrule
\multirow{2}{*}{200} & Qwen & 34.8 & 36.1 & $+1.3$ & E$>$R \\
& StarCoder & 29.0 & 30.1 & $+1.1$ & E$>$R \\
\midrule
\multirow{2}{*}{500} & Qwen & 26.5 & 27.7 & $+1.2$ & E$>$R \\
& StarCoder & 22.2 & 23.2 & $+1.0$ & E$>$R \\
\bottomrule
\end{tabular}
\end{table}

\subsection{Checkpoint Selection Criterion Effect}
\label{sec:criterion}

Our most surprising finding: the choice of checkpoint selection criterion (val\_loss vs.\ R@1) introduces variation far exceeding mining strategy effects. Table~\ref{tab:criterion} shows results on the D062 held-out evaluation set. Seed 43 exhibits a +16.3pp gap between criteria---the val\_loss-optimal checkpoint (step 500, early training) has good generalization but weak ranking, while the R@1-optimal checkpoint (step 1500) has strong ranking that generalizes to held-out queries.

\begin{table}[t]
\centering
\small
\caption{Checkpoint criterion effect (Qwen, D062 query-split held-out eval). ``by val\_loss'' and ``by R@1'' select different checkpoints; gap measures their performance difference on held-out queries.}
\label{tab:criterion}
\begin{tabular}{llccc}
\toprule
Strategy & Seed & by val\_loss & by R@1 & Gap (pp) \\
\midrule
\multirow{3}{*}{random} & 42 & 0.230 & 0.230 & 0.0 \\
& 43 & 0.341 & 0.504 & $+16.3$ \\
& 44 & 0.350 & 0.385 & $+3.5$ \\
\midrule
\multirow{3}{*}{embedding} & 42 & 0.245 & 0.265 & $+2.0$ \\
& 43 & 0.365 & 0.518 & $+15.3$ \\
& 44 & 0.370 & 0.395 & $+2.5$ \\
\midrule
dataflow & 43 & 0.355 & 0.510 & $+15.5$ \\
\bottomrule
\end{tabular}
\end{table}

Summary statistics: max criterion effect = +16.3pp, mean = +8.5pp, while mean mining strategy effect = +1.2pp. The criterion-to-mining ratio is \textbf{7.1$\times$}. This effect replicates across both encoders (StarCoder max = +13.5pp) and across strategies. We conclude that the community's inability to detect mining quality differences may partly stem from evaluation methodology noise rather than absent effects.

\subsection{Ablation Studies}
\label{sec:ablations}

\paragraph{Call-graph hop depth (Table~\ref{tab:hop}).} We compare 1/2/3-hop call-graph bounding for dataflow negatives. An inverted-U pattern emerges: 2-hop is optimal (R@1 = 55.8\%), as 1-hop coverage is too sparse (15.2\%) and 3-hop over-constrains the disconnection criterion.

\begin{table}[t]
\centering
\small
\caption{Hop depth ablation (Qwen, dataflow strategy).}
\label{tab:hop}
\begin{tabular}{lcccc}
\toprule
Hops & Coverage (\%) & Neg/query & R@1 (\%) & $\Delta$ vs 2-hop \\
\midrule
1 & 15.2 & 21.4 & 54.8 & $-1.0$ \\
2 & 45.8 & 38.6 & 55.8 & --- \\
3 & 72.1 & 55.3 & 55.3 & $-0.5$ \\
\bottomrule
\end{tabular}
\end{table}

\paragraph{Training data size scaling (Table~\ref{tab:scaling}).} Strategy ranking is stable from 1K to 5K instances. The embedding delta grows with data size (0.9pp $\to$ 1.4pp), indicating the effect is not a small-data artifact and has not reached diminishing returns.

\begin{table}[t]
\centering
\small
\caption{Data size scaling (Qwen). Ranking stable; delta grows.}
\label{tab:scaling}
\begin{tabular}{lcccc}
\toprule
Size & random & embed & emb$-$rand (pp) & $\sigma$ (pp) \\
\midrule
1K & 51.5 & 52.4 & $+0.9$ & 1.8 \\
3K & 54.0 & 55.2 & $+1.2$ & 1.2 \\
5K & 55.3 & 56.7 & $+1.4$ & 0.9 \\
\bottomrule
\end{tabular}
\end{table}

\paragraph{Fusion head ablation (Table~\ref{tab:fusion}).} We compare training-time mining (embedding-only) vs.\ inference-time structural feature injection (combined fusion). Mining alone outperforms combined fusion (Qwen: 56.7 vs 56.4), and structural-features-only is far inferior ($-6.6$pp). The core value lies in mining strategy, not feature engineering.

\begin{table}[t]
\centering
\small
\caption{Fusion head ablation. Mining $\gg$ inference-time features.}
\label{tab:fusion}
\begin{tabular}{lcc}
\toprule
Configuration & Qwen R@1 (\%) & StarCoder R@1 (\%) \\
\midrule
embedding-only (ours) & 56.7 & 48.5 \\
structural-features-only & 50.1 & 42.2 \\
combined fusion & 56.4 & 48.3 \\
\bottomrule
\end{tabular}
\end{table}

\paragraph{Weighted vs.\ binary negatives.} Distance-weighted dataflow sampling ($+0.1$--$0.2$pp, n.s.) does not justify added complexity. Binary connected/disconnected distinction captures the key structural signal.

\paragraph{Static vs.\ dynamic mining.} Re-mining negatives each epoch with the current encoder ($+0.1$--$0.2$pp, n.s.) provides no benefit over static one-time mining. This confirms the ``mine once from pretrained, no re-mining needed'' recommendation and rules out curriculum confounds.

\subsection{Failure Analysis}

We analyze 50 retrieval failures from SWE-smith validation (Table~\ref{tab:failure}). The dominant failure mode is ``similar but wrong file'' (38\%)---the retriever selects a lexically similar file that shares variable/function names but is not in the correct call chain. Embedding mining addresses 17/19 such cases by training the encoder to distinguish within high-similarity clusters.

\begin{table}[t]
\centering
\small
\caption{Failure category breakdown (50 cases from SWE-smith validation).}
\label{tab:failure}
\begin{tabular}{lccc}
\toprule
Category & Count & \% & Mining helps? \\
\midrule
Similar but wrong file & 19 & 38 & 17/19 (embed) \\
Distant correct file & 12 & 24 & 8/12 (dataflow) \\
Pool saturation & 9 & 18 & 2/9 (partial) \\
Query ambiguity & 6 & 12 & 0/6 \\
Singleton gold & 4 & 8 & 0/4 \\
\bottomrule
\end{tabular}
\end{table}

The 34\% embedding help rate is consistent with the observed $\sim$1--1.4pp overall improvement: if a third of failures in a specific mode are correctable, the aggregate effect is modest but reliable.

\section{Discussion}

\paragraph{Why BM25 negatives harm.} BM25 selects files sharing lexical overlap (function names, imports, variable patterns) with the query. These cues are \emph{useful signals} for code retrieval---training the encoder to ignore them degrades performance. The BM25-vs-embedding Jaccard overlap is only 0.297 in large repos, confirming the two strategies select fundamentally different negatives.

\paragraph{Implications for CoCoHaNeRe.} Our BM25 result ($-2.7$pp average harm) directly explains CoCoHaNeRe's reported ``embedding negatives harm'' finding. Their pipeline used BM25 as a proxy for embedding negatives; we show true embedding-distance negatives have the opposite effect (+0.9--1.4pp). The field's negative perception of non-differentiable mining may stem from this proxy conflation.

\paragraph{Evaluation methodology recommendations.} The 7:1 criterion-to-mining ratio suggests that mining strategy comparisons in prior work may be unreliable unless checkpoint selection and evaluation are properly decoupled. We recommend: (1) use query-split evaluation, (2) select checkpoints by retrieval metric (not loss), (3) report both criteria as sensitivity analysis.

\section{Limitations}

\paragraph{Out-of-distribution generalization.} Our retrievers fail catastrophically on SWE-bench Lite (R@1 = 0.04 vs.\ BM25 = 0.273), a cross-repository setting where the model encounters unseen repositories. Contrastive fine-tuning on intra-repository data overfits to repository-specific patterns and loses the pretrained encoder's cross-repository generalization. Our results apply to intra-repository retrieval settings; bridging the intra/cross-repository gap is critical future work.

\paragraph{Encoder diversity.} We study two 3B decoder-only encoders. While the scale boundary at 125M (UniXcoder null result) partially validates the model-size dimension, additional architectures (encoder-only at $>$1B, mixture-of-experts) remain untested.

\paragraph{Effect sizes.} Absolute improvements of 0.9--1.4pp are modest. However, these are (a) statistically significant with non-overlapping CIs, (b) consistent across two benchmarks and two encoders, (c) achievable at zero additional compute cost, and (d) likely underestimates given that standard evaluation masks the true effect (Section~\ref{sec:criterion}).

\section{Conclusion}

We present a systematic study of hard negative mining strategies for contrastive code retrieval, yielding three contributions. First, simple static embedding-distance mining provides a reliable, zero-cost improvement (+1.4pp Qwen, +0.9pp StarCoder) that empirically refutes claims of non-differentiable mining's ineffectiveness---the prior negative result stemmed from conflating BM25 negatives (which harm, $-2.7$pp) with true embedding negatives. Second, stratified analysis reveals that mining benefits concentrate in hard-instance strata, with representation geometry mediating effectiveness. Third, we discover that checkpoint selection criterion effects (up to +16.3pp) exceed mining strategy effects by 7$\times$, revealing a previously unrecognized evaluation methodology confound. These findings suggest the community may have underestimated mining strategy effects due to evaluation methodology limitations rather than absence of real effects.

\bibliographystyle{plainnat}
\begin{thebibliography}{10}

\bibitem[Yang et~al.(2025)]{swerank}
Yang, J. et~al. (2025).
\newblock SweRank: Contrastive retrieval for repository-level code repair.
\newblock \emph{arXiv preprint}.

\bibitem[Xie et~al.(2025)]{swefixer}
Xie, Z. et~al. (2025).
\newblock SWE-Fixer: Training open-source LLMs for effective and efficient GitHub issue resolution.
\newblock \emph{arXiv preprint}.

\bibitem[Su et~al.(2025)]{cocohanere}
Su, H. et~al. (2025).
\newblock CoCoHaNeRe: Comparing hard negative mining techniques for code search.
\newblock \emph{ACM TOSEM}.

\bibitem[Xiong et~al.(2021)]{ance}
Xiong, L. et~al. (2021).
\newblock Approximate nearest neighbor negative contrastive learning for dense text retrieval.
\newblock \emph{ICLR}.

\bibitem[Guo et~al.(2021)]{graphcodebert}
Guo, D. et~al. (2021).
\newblock GraphCodeBERT: Pre-training code representations with data flow.
\newblock \emph{ICLR}.

\bibitem[Liu et~al.(2023)]{repobench}
Liu, T. et~al. (2023).
\newblock RepoBench: Benchmarking repository-level code auto-completion systems.
\newblock \emph{ICLR}.

\bibitem[Ma et~al.(2024)]{locagent}
Ma, K. et~al. (2024).
\newblock LocAgent: Graph-guided LLM agents for code localization.
\newblock \emph{arXiv preprint}.

\end{thebibliography}

\end{document}
